\chapter{Message order and the security it gives}
\label{chap:variants}

What can be proved about the protocol of \cite{JKO13} depends less on the garbling scheme than on
the order of the messages: when the prover's witness bits become fixed, compared with the moment
it sees the garbled circuit, and when the verifier opens its garbling, compared with the moment
the prover discloses its output label. The literature fixes these choices in different places, often
only in a figure or a footnote. This chapter collects them, says which kind of security each
choice allows, and points to the statements of this library that cover it. The chapter is
expository and is not itself formalized. It contains no new statement; every claim about the
library refers to a theorem proved in a later chapter, and a phrase such as ``soundness needs
adaptive authenticity'' describes what the proofs here use, not a proved necessity.

\section{The protocol in order}

In the protocol of Jawurek, Kerschbaum and Orlandi \cite[Sect.~1.1, Fig.~1]{JKO13} the messages
are, in order:
\begin{enumerate}
  \item the prover runs the oblivious transfers as receiver, one per bit of its witness $w$; the
  verifier, as sender, inputs the two labels of each input wire;
  \item the verifier sends the garbled circuit (with a two-round OT, together with its second
  message);
  \item the prover evaluates and \emph{commits} to the output label it obtained;
  \item the verifier \emph{opens} its garbling: the seed, and all its inputs to the oblivious
  transfers (``open-all'');
  \item the prover checks that the opened garbling is a garbling of the agreed circuit and agrees
  with the labels it received, and only then \emph{reveals} its committed label;
  \item the verifier accepts if the label is the output label for the value $1$.
\end{enumerate}
Three choices determine the security: the position of step~2 with respect to the moment the
witness bits are fixed in step~1; whether step~4 exists and comes after step~3; and how the
commitment and the oblivious transfer are realized.

\section{When the witness bits are fixed: static or adaptive authenticity}

\emph{Witness fixed before the garbled circuit is seen.} If the receiver's message in the OT
determines its choice bits, and the garbled circuit is sent only afterwards, the prover's input
does not depend on the garbling. Soundness then needs only the static authenticity of the garbling:
the input is chosen first, and the adversary must forge the other output label. This is the order of
\cite{JKO13}, who note that the garbled circuit may travel with the second message of a two-round
OT \cite[Sect.~1.1]{JKO13}, and that sending it before the oblivious transfers would require a
garbling secure against adaptive attacks \cite[Sect.~3, footnote~7]{JKO13}. In this library: soundness reduces exactly to
the authenticity game (Theorem~\ref{thm:sound}), bounded by $(q+1)/(|V|-1)$ in the random oracle
model (Theorems~\ref{thm:rom-bound} and~\ref{thm:garble-rom}). The PVW OT in messy mode is of this
kind: the receiver's key determines its choice, which a trapdoor extracts (Lemma~\ref{lem:pvw-ext},
Theorem~\ref{thm:pvw-sound}).

\emph{Witness fixed after the garbled circuit is seen.} If the garbled circuit is sent before the
oblivious transfers, or if the OT lets the receiver postpone its choice, the prover may choose its
witness bits as a function of the garbled circuit. Soundness then needs \emph{adaptive}
authenticity, in which the adversary sees the garbled circuit and requests its input labels one at a
time \cite[Sect.~4, Fig.~6]{BHR12}. The Chou--Orlandi OT is of this kind: the receiver's message does
not fix its choice, and a receiver can delay the decryption until after it has seen the garbled
circuit (\cite[Sect.~3]{BPRS17}, \cite[App.~A]{DKLs18}; see also \cite[Sect.~1.1]{CO15}). In this library: adaptive authenticity
for the garbling formalized here, with bound $(q+1)/(|V|-1) + q/|V|$ for $q$ actions
(Theorem~\ref{thm:adaptive}), and soundness with the Chou--Orlandi OT, by reduction to square
Diffie--Hellman (Theorem~\ref{thm:co-sound}).

\section{When the garbling is opened: honest or malicious verifier}

\emph{No opening.} If the prover sends its output label directly and the verifier never opens the
garbling, the protocol is zero knowledge only against an honest verifier \cite[Sect.~3]{JKO13}: a
malicious verifier can garble a circuit other than the agreed one, or mount a selective-failure
attack in the input phase \cite[Sect.~1]{JKO13}. In this library: honest-verifier zero knowledge
(Theorem~\ref{thm:hvzk}), and an explicit malicious verifier that learns a witness bit from the OT
pairs alone (Proposition~\ref{prop:no-check-leaks}).

\emph{Commit, open, check, reveal.} With steps~3 to~5 in this order the protocol is zero knowledge
against a malicious verifier, with a single garbled circuit. The opening is possible because in zero
knowledge the verifier has no secret input, unlike in general two-party computation
\cite[Sect.~1]{JKO13}. The order matters: the prover must commit to its label before the opening,
check the opening, and reveal the label only afterwards. The OT must allow this opening, so that the
prover can check that the labels it received are those of the opened garbling
\cite[Sect.~1.1]{JKO13}. In this library: perfect zero knowledge in the hybrid model
(Theorem~\ref{thm:zk}).

\emph{The garbling cannot be opened.} When the garbler has inputs or keys that must stay secret, the
garbling cannot be opened at the end, and a malicious garbler is handled by cut-and-choose: several
garbled circuits, some opened and some evaluated. If one good evaluated circuit suffices, choosing
each circuit with an independent fair coin gives error $2^{-N}$ \cite{Lin13}. In this library:
Theorems~\ref{thm:cc-fixed}, \ref{thm:cc-lindell} and~\ref{thm:caught-or-correct}.

\section{How the commitment and the oblivious transfer are realized}

\emph{Ideal functionalities.} With ideal OT and ideal commitment, zero knowledge is perfect and
soundness is exactly the authenticity game (Theorems~\ref{thm:zk} and~\ref{thm:sound}).

\emph{Hash commitment in the random oracle model.} Zero knowledge becomes statistical, with error
at most $(q_1+q_2+q_3)/|R|$ for a verifier making $q_1+q_2+q_3$ oracle queries; the simulator
\emph{programs} the oracle (Theorem~\ref{thm:zk-rom}). Soundness goes through an extractor that
reads the prover's queries, at an extra cost $(q_1^2+q_2+1)/|D|$ for $q_1$ queries before the
commitment and $q_2$ after it ($n_1$ and $n_2$ in Theorem~\ref{thm:e2e}) (Theorems~\ref{thm:binding}
and~\ref{thm:e2e}).

\emph{PVW OT with a common string.} In messy mode the receiver's choice is extractable, and
soundness keeps the bound of the hybrid model (Theorem~\ref{thm:pvw-sound}). Zero knowledge is
computational: it is reduced to DDH by two explicit distinguishers between the modes
(Theorems~\ref{thm:pvw-zk-modes} and~\ref{thm:pvw-zk-ddh}). The receiver key $(0, 0)$ sends the messages in the
clear, and the sender in this library rejects every key whose first component is $0$
(Lemma~\ref{lem:pvw-ext}).

\emph{Chou--Orlandi OT in the random oracle model, without common string.} Receiver privacy is
perfect and zero knowledge holds with the real OT and the hash commitment
(Theorem~\ref{thm:zk-ot}). Soundness needs adaptive authenticity, and costs $k/q$ for a secret
$a_j = 0$ plus the advantages of explicit square Diffie--Hellman solvers
(Theorem~\ref{thm:co-sound}). Chou and Orlandi withdraw the claim of UC security of their protocol
\cite[Sect.~1.1]{CO15}.

\section{When parties are corrupted}
\label{sec:corruptions}

The security theorems of this library assume \emph{static} corruptions: the adversary decides at
the start which party it controls. Chapter~\ref{chap:corruption} proves some ingredients for
adaptive corruptions, but no security statement against them. In
the framework of universal composability (UC) party
corruption is modelled by conventions on top of the basic model of execution
\cite[Sect.~7.1]{Can20}. The formulation given there captures \emph{adaptive} corruptions: the
adversary corrupts parties as the computation proceeds, based on what it has seen, and receives the
entire local state of a corrupted party. Static corruptions, where the corrupted parties are fixed
before the computation starts, are a weaker threat model, and the setting where erasures are not
trusted corresponds to protocols that write each data cell at most once
\cite[Sect.~7.1.1]{Can20}. Against adaptive corruptions a simulator must produce the transcript
without the witness and, when the prover is corrupted later, \emph{explain} it: produce internal
randomness consistent with the actual witness \cite[Sect.~1.2]{GKPS18}.

\emph{Why it matters with two parties.} We add a remark of our own. In a single proof between two
parties the verifier has no secret input, and a corrupted prover hands over its witness anyway.
Adaptive security matters most when the proof is a component of a larger protocol, with many
sessions or many parties, in which devices may be compromised at different times: the security
proof of the larger protocol, through the composition theorem \cite[Theorem~22]{Can20}, then relies
on every component being simulatable under late corruption. Ganesh, Kondi, Patra and Sarkar
motivate adaptive security by adversaries that use their resources gradually, and by applications
that use the protocol of \cite{JKO13} as a component \cite[Sect.~1.3]{GKPS18}.

\emph{What has to be explained, and when.} In the protocol of \cite{JKO13} the verifier has no
input, so its state never needs to be faked, and once the prover holds its output label nothing it
does depends on the witness \cite[Sect.~3.2]{GKPS18}; in this library the revealed label is the
simulator's output (Theorem~\ref{thm:zk}). What must be explained therefore depends on the moment
the prover is corrupted. The table is our reading of the protocol; \cite{GKPS18} prove the case of
an ideal commitment, stated after it.
\begin{center}
\begin{tabular}{|p{0.3\textwidth}|p{0.3\textwidth}|p{0.3\textwidth}|}
  \hline
  \textbf{Prover corrupted} & \textbf{To be explained} & \textbf{Property used} \\
  \hline
  before the oblivious transfers & nothing & none \\
  \hline
  after the oblivious transfers, before the commitment & its OT messages, for the actual witness &
  receiver-equivocal OT \\
  \hline
  after the commitment, before the opening & its OT messages and its commitment & also an equivocal
  commitment \\
  \hline
  after the label is revealed & its OT messages; the commitment is already opened &
  receiver-equivocal OT; the label does not depend on the witness \\
  \hline
\end{tabular}
\end{center}
A receiver-equivocal OT is a secure OT whose receiver message can be produced, from the trapdoor of
the setup, without the choice bit and later opened to either bit
\cite[Sect.~3, Definition~3.1]{GKPS18}; the PVW OT in decryption mode is one, assuming DDH
\cite[Sect.~3.3, Theorem~3.2]{GKPS18}. DDH is used for the security of the OT; the equivocation
itself is exact (Theorem~\ref{thm:pvw-re}). With an ideal commitment there is nothing
to explain. A hash commitment in the \emph{programmable} random oracle model can be equivocated by
setting the oracle at the committed point after the fact, a technique recalled in
\cite[Sect.~1.2]{GKPS18}. With a receiver-equivocal OT and an ideal commitment, the
protocol of \cite{JKO13} is secure against adaptive corruptions \cite[Sect.~3]{GKPS18}; after
the execution, the only work is to equivocate the prover's view of the OT
\cite[Sect.~3.2]{GKPS18}.

\emph{The other choices.} The order of the messages interacts with the corruption time. For the
Chou--Orlandi OT we note the following, which follows from the form of the messages: the receiver
message $B = b \cdot g + wA$, with $A = a \cdot g$, is also $(b + (2w - 1) a) \cdot g + (1 - w) A$, so
explaining it for the other choice bit requires the sender's secret $a$. The simulator knows $a$ when
the verifier is honest, and otherwise learns it only at the open-all step, so a corruption of the
prover before the opening is the delicate case. If the garbler has an input of its own, as in general
two-party computation, the garbled circuit itself must be equivocated, which requires much heavier
tools \cite[Sect.~3.2]{GKPS18}. Without erasures there is more to explain, and adaptively secure
oblivious transfer is harder to build \cite{BPRS17}.

\emph{In this library.} Adaptive security itself is not formalized: neither the UC framework, nor
a model of execution with corruptions. Chapter~\ref{chap:corruption} proves the protocol-specific
facts behind the table: the receiver of the PVW OT in decryption mode, with the trapdoor of the
setup, and the receiver of the Chou--Orlandi OT, with the sender's secret, satisfy the
equivocation condition on messages and randomness (Theorem~\ref{thm:pvw-re},
Proposition~\ref{prop:co-re}); under that condition the prover's OT messages, and the verifier's
view of the prover, can be explained for any witness (Theorem~\ref{thm:explain-prover}); the hash commitment can
be equivocated by programming the oracle (Theorem~\ref{thm:equivocate-rom}); and soundness holds
with the PVW OT in decryption mode, up to the DDH advantage (Theorem~\ref{thm:pvw-dec-sound}).
These are ingredients only; they are not composed into a simulation. A somewhat restricted variant of the UC framework has been
mechanized in EasyCrypt by Canetti, Stoughton and Varia, with a worked example \cite{CSV19}. The
README of the public core of CatCrypt lists machinery for universal composability in Lean, and says
that the universal-composability theory itself is in separate distributions \cite{CCcore}.

\section{Summary}

\begin{center}
\begin{tabular}{|p{0.27\textwidth}|p{0.3\textwidth}|p{0.33\textwidth}|}
  \hline
  \textbf{Choice} & \textbf{Security that can be expected} & \textbf{In this library} \\
  \hline
  Witness fixed, then garbled circuit & Soundness from static authenticity & Theorems~\ref{thm:sound}, \ref{thm:garble-rom}, \ref{thm:e2e}, \ref{thm:pvw-sound} \\
  \hline
  Garbled circuit first, or OT with late choice & Soundness needs adaptive authenticity & Theorems~\ref{thm:adaptive}, \ref{thm:co-sound} \\
  \hline
  No opening & Honest-verifier zero knowledge only & Theorem~\ref{thm:hvzk}, Proposition~\ref{prop:no-check-leaks} \\
  \hline
  Commit, open, check, reveal & Zero knowledge against a malicious verifier, one circuit & Theorems~\ref{thm:zk}, \ref{thm:zk-rom}, \ref{thm:zk-ot}, \ref{thm:pvw-zk-ddh} \\
  \hline
  Garbling cannot be opened & Cut-and-choose, error $2^{-N}$ & Theorems~\ref{thm:cc-lindell}, \ref{thm:caught-or-correct} \\
  \hline
  Ideal OT and commitment & Perfect zero knowledge & Theorem~\ref{thm:zk} \\
  \hline
  Hash commitment (random oracle) & Statistical zero knowledge, programmable oracle & Theorem~\ref{thm:zk-rom} \\
  \hline
  PVW OT, messy mode & Computational zero knowledge (DDH) & Theorem~\ref{thm:pvw-zk-ddh} \\
  \hline
  Chou--Orlandi OT & Soundness by reduction to square Diffie--Hellman & Theorem~\ref{thm:co-sound} \\
  \hline
  Adaptive corruptions & A receiver-equivocal OT with an ideal commitment suffices \cite{GKPS18} & Some ingredients:
  Theorems~\ref{thm:pvw-re}, \ref{thm:explain-prover}, \ref{thm:equivocate-rom}, \ref{thm:pvw-dec-sound}; adaptive
  security itself is not formalized (Section~\ref{sec:corruptions}) \\
  \hline
\end{tabular}
\end{center}

Security against adaptive corruptions, a different axis, is discussed in
Section~\ref{sec:corruptions}.
